% Chapter 8: Antithetic variates, moment matching and control variates.
\providecommand{\Corr}{\operatorname{Corr}}

The plain Monte Carlo estimator of \cref{ch:06-monte-carlo} has standard error
$\sigma/\sqrt n$: one more correct digit costs a hundred times more samples. Independent
sampling cannot beat the rate $n^{-1/2}$, but it can beat the constant. This chapter is
about making $\sigma^2$, or rather the product of $\sigma^2$ and the cost per sample,
smaller by using what we know about the problem.

Four ideas, each simple to state. \emph{Antithetic variates} use the symmetry of the
normal law: pair every draw $Z$ with $-Z$. \emph{Moment matching} forces the sample to
have the moments of the model. \emph{Common random numbers} correlate the estimators of
two related prices, so that their difference is accurate even when each price is not.
\emph{Control variates} subtract a correlated quantity whose mean is known exactly, and
include, as limiting cases, regression on several such quantities and hedging. Each comes
with an exact variance formula, a condition under which it helps, and examples in which it
does not. They combine, but not multiplicatively; \cref{sec:cv-lab} measures all of them
on the course's test case.

\section{Setting and running examples}\label{sec:cv-setting}

We estimate $\mu=\E[f(Z)]$, where $Z$ is a standard normal vector in $\R^d$ (one
coordinate per time step or per asset) and $f(Z)\in L^2$ with $\sigma^2=\Var f(Z)>0$. The
plain estimator $\hat\mu_n=\frac1n\sum_{i=1}^nf(Z_i)$ of \cref{def:mc-estimator} has
variance $\sigma^2/n$. Every method below produces another unbiased (or asymptotically
unbiased) estimator. We report its \emph{standard-error ratio} to $\hat\mu_n$ at equal
$n$, $\varrho=\sigma_2/\sigma$, and, if one of its samples costs $c$ times a plain one, its
efficiency relative to plain Monte Carlo (\cref{def:mc-efficiency})
\begin{equation}\label{eq:cv-efficiency}
  \mathcal E_{2,1}=\frac{\sigma^2\tau_1}{\sigma_2^2\tau_2}=\frac{1}{c\,\varrho^2}.
\end{equation}
Following the course, we take the conservative value $c=2$ for every method when a cost
must be assumed; the notebook measures it.

Throughout, the test case is $S=K=100$, $T=1$, $r=1\%$, $y=2\%$, $\sigma=20\%$, and
\[
  S_T=S\exp\big(\nu T+s Z\big),\qquad \nu=r-y-\tfrac{\sigma^2}{2}=-0.03,\qquad s=\sigma\sqrt T=0.2,
\]
with $Z\sim\mathcal N(0,1)$ (\cref{prop:ps-gbm-exact}). The running examples are
\begin{enumerate}[label=(\alph*)]
  \item the \emph{stock}, $f(Z)=e^{-rT}S_T$, with $\mu=Se^{-yT}=98.0199$: trivial to
    price, and therefore a good test;
  \item the \emph{call}, $f(Z)=e^{-rT}(S_T-K)^+$, with $\mu=7.36429$ and
    $\sigma^2=158.006$ (for the undiscounted payoff, $161.198$), so $\sigma=12.570$;
  \item the \emph{square-root call} of \cref{eg:cv-sqrt}.
\end{enumerate}

\section{Antithetic variates}\label{sec:cv-antithetic}

If $Z\sim\mathcal N(0,I_d)$ then $-Z$ has the same law. Draw $Z_1,\dots,Z_n$ independently
and use each draw twice:
\begin{equation}\label{eq:cv-antithetic}
  \hat\mu^{\mathrm a}_n=\frac1n\sum_{i=1}^n\frac{f(Z_i)+f(-Z_i)}{2}.
\end{equation}
For a path driven by the increments $Z=(Z_1,\dots,Z_m)$, the antithetic path is driven by
$-Z$: it is the original path reflected in its drift. With uniforms, $U$ is paired with
$1-U$, which under inversion (\cref{prop:rng-inversion}) is the same as $Z$ with $-Z$. The
estimator averages $n$ i.i.d.\ \emph{pair averages}, so its standard error comes from the
$n$ pair averages, never from the $2n$ individual payoffs, which are dependent.

When does pairing help? The answer rests on an inequality that goes back to Chebyshev in
one dimension, and in several dimensions is the simplest case of \emph{association}
\citep{esary-proschan-walkup}.

\begin{lemma}[Monotone functions are positively correlated]\label{lem:cv-association}
Let $Z=(Z_1,\dots,Z_d)$ have independent components, and let $g,h:\R^d\to\R$ be Borel
functions, nondecreasing in each coordinate, with $g(Z),h(Z)\in L^2$. Then $\Cov(g(Z),h(Z))\ge0$.
\end{lemma}

\begin{proof}
Assume first that $g$ and $h$ are bounded, and argue by induction on $d$. For $d=1$, let
$Z'$ be an independent copy of $Z$. Since $g$ and $h$ are nondecreasing, $g(Z)-g(Z')$ and
$h(Z)-h(Z')$ have the same sign, so $(g(Z)-g(Z'))(h(Z)-h(Z'))\ge0$; taking expectations,
$2\E[g(Z)h(Z)]-2\E[g(Z)]\E[h(Z)]\ge0$.

For $d\ge2$, write $Z=(W,Z_d)$ with $W=(Z_1,\dots,Z_{d-1})$ and, for $w\in\R^{d-1}$,
\[
  G(w)=\E[g(w,Z_d)],\qquad H(w)=\E[h(w,Z_d)],\qquad Q(w)=\E[g(w,Z_d)h(w,Z_d)] .
\]
These are bounded and measurable (Fubini), and $G$, $H$ are nondecreasing in each
coordinate of $w$ (integrate the pointwise inequalities). For fixed $w$, the functions
$t\mapsto g(w,t)$ and $t\mapsto h(w,t)$ are nondecreasing, so the case $d=1$ gives
$Q(w)\ge G(w)H(w)$. Since $W$ and $Z_d$ are independent, Fubini's theorem gives
$\E[g(Z)h(Z)]=\E[Q(W)]$, $\E[g(Z)]=\E[G(W)]$ and $\E[h(Z)]=\E[H(W)]$. Hence, by the
induction hypothesis applied to $G$ and $H$,
\[
  \E[g(Z)h(Z)]=\E[Q(W)]\ge\E[G(W)H(W)]\ge\E[G(W)]\,\E[H(W)]=\E[g(Z)]\,\E[h(Z)] .
\]
For general $g,h$, the truncations $g_k=\max(\min(g,k),-k)$ and $h_k$ are bounded and
still nondecreasing, so $\Cov(g_k(Z),h_k(Z))\ge0$; as $k\to\infty$, dominated convergence
($|g_kh_k|\le|gh|$, $|g_k|\le|g|$) gives the claim.
\end{proof}

\begin{proposition}[Antithetic variates]\label{prop:cv-antithetic}
Let $Z\eqd-Z$, $f(Z)\in L^2$ with $\sigma^2=\Var f(Z)>0$, and
$\rho=\Corr\big(f(Z),f(-Z)\big)$. Write $f=f_{\mathrm e}+f_{\mathrm o}$ with the even and
odd parts $f_{\mathrm e}(z)=\frac12(f(z)+f(-z))$ and $f_{\mathrm o}(z)=\frac12(f(z)-f(-z))$.
\begin{enumerate}[label=(\roman*)]
  \item $f_{\mathrm e}(Z)$ and $f_{\mathrm o}(Z)$ are uncorrelated, so
    $\sigma^2=\Var f_{\mathrm e}(Z)+\Var f_{\mathrm o}(Z)$, and
    \[
      \rho=\frac{\Var f_{\mathrm e}(Z)-\Var f_{\mathrm o}(Z)}{\Var f_{\mathrm e}(Z)+\Var f_{\mathrm o}(Z)} .
    \]
  \item $\hat\mu^{\mathrm a}_n$ is unbiased and
    \begin{equation}\label{eq:cv-antithetic-var}
      \Var\big(\hat\mu^{\mathrm a}_n\big)=\frac{\Var f_{\mathrm e}(Z)}{n}=\frac{\sigma^2}{n}\,\frac{1+\rho}{2} .
    \end{equation}
  \item If the components of $Z$ are independent and $f$ is monotone in each coordinate
    (nondecreasing in some, nonincreasing in the others), then $\rho\le0$: $n$ antithetic
    pairs have at most the variance of $2n$ independent draws.
\end{enumerate}
\end{proposition}

\begin{proof}
(i) Since $f(-Z)\eqd f(Z)$, $\E f_{\mathrm o}(Z)=0$ and
$\E[f_{\mathrm e}(Z)f_{\mathrm o}(Z)]=\frac14\E\big[f(Z)^2-f(-Z)^2\big]=0$. So
$\sigma^2=\Var(f_{\mathrm e}+f_{\mathrm o})(Z)$ splits into the two variances, and
$\Cov(f(Z),f(-Z))=\Cov\big((f_{\mathrm e}+f_{\mathrm o})(Z),(f_{\mathrm e}-f_{\mathrm o})(Z)\big)
=\Var f_{\mathrm e}(Z)-\Var f_{\mathrm o}(Z)$. Divide by $\sigma^2=\Var f(-Z)$.

(ii) The pair averages $f_{\mathrm e}(Z_i)$ are i.i.d.\ with mean $\mu$. By (i),
$\Var f_{\mathrm e}+\Var f_{\mathrm o}=\sigma^2$ and $\Var f_{\mathrm e}-\Var f_{\mathrm o}=\rho\sigma^2$, so
$\Var f_{\mathrm e}(Z)=\sigma^2(1+\rho)/2$.

(iii) Choose signs $\varepsilon_j\in\{\pm1\}$ such that
$\tilde f(v)=f(\varepsilon_1v_1,\dots,\varepsilon_dv_d)$ is nondecreasing in every coordinate, and
let $V=(\varepsilon_jZ_j)_j$, which again has independent components. Then $f(Z)=\tilde f(V)$
and $f(-Z)=\tilde f(-V)$. Both $\tilde f$ and $v\mapsto-\tilde f(-v)$ are nondecreasing in every
coordinate, so \cref{lem:cv-association} gives $\Cov(\tilde f(V),-\tilde f(-V))\ge0$, that is
$\rho\le0$.
\end{proof}

The even--odd split says exactly what antithetic sampling does: it integrates the odd part
of $f$ with no error at all, and samples the even part as before. It helps when $f$ is
nearly odd about the centre of the distribution (nearly linear payoffs are the main
example), and it is useless
when $f$ is even.

\paragraph{Cost and efficiency.} A pair costs $c$ times a single draw with $1\le c\le2$:
$c\approx2$ when evaluating the payoff dominates, and less when generating normals or paths
dominates, since the antithetic path reuses them. By \eqref{eq:cv-efficiency} and
\eqref{eq:cv-antithetic-var}, the efficiency of $n$ pairs against $n$ plain draws is
$2/(c(1+\rho))$; at $c=2$ it is $1/(1+\rho)$, which is at least $1$ under the hypotheses of
(iii).

\begin{example}[The stock]\label{eg:cv-stock}
For $f(z)=e^{-rT}Se^{\nu T}e^{sz}$ the even and odd parts are multiples of $\cosh(sz)$ and
$\sinh(sz)$. With $\E\cosh(sZ)=e^{s^2/2}$, $\E\cosh^2(sZ)=\frac12(e^{2s^2}+1)$ and
$\E\sinh^2(sZ)=\frac12(e^{2s^2}-1)$,
\[
  \Var\cosh(sZ)=\tfrac12(e^{s^2}-1)^2,\qquad \Var\sinh(sZ)=\tfrac12(e^{2s^2}-1),
  \qquad\text{so}\qquad \rho=-e^{-s^2} .
\]
For $s=0.2$: $\rho=-0.9608$, the standard-error ratio is
$\sqrt{(1+\rho)/2}=\sqrt{(1-e^{-s^2})/2}=0.1400\approx s/\sqrt2$, and the efficiency at
$c=2$ is $1/(1-e^{-s^2})=25.5$. These are the course's ``about 14\%'' and ``about 25'':
over the range of $sZ$ the exponential is nearly linear, and linear functions are
integrated exactly.
\end{example}

For payoffs that vanish on part of the line, the antithetic correlation has a closed form.
This is the course's exam formula.

\begin{proposition}[Disjoint payoffs]\label{prop:cv-disjoint}
In the setting of \cref{prop:cv-antithetic}, suppose that $f(Z)f(-Z)=0$ almost surely, and
let $\mu=\E f(Z)$. Then $\mu^2\le\sigma^2$ and
\[
  \rho=-\frac{\mu^2}{\sigma^2},\qquad
  \Var\big(\hat\mu^{\mathrm a}_n\big)=\frac{\sigma^2-\mu^2}{2n},\qquad
  \frac{\sqrt{\Var\hat\mu^{\mathrm a}_n}}{\sqrt{\Var\hat\mu_n}}=\frac{1}{\sqrt2}\sqrt{1-\frac{\mu^2}{\sigma^2}},
\]
and the efficiency at $c=2$ is $\sigma^2/(\sigma^2-\mu^2)\ge1$ (infinite if
$\mu^2=\sigma^2$). In the Black--Scholes model the call $(S_T-K)^+$ and the digital call
$\1_{\{S_T>K\}}$ satisfy the hypothesis if and only if $K\ge Se^{\nu T}$, the median of
$S_T$; the put and the digital put if and only if $K\le Se^{\nu T}$.
\end{proposition}

\begin{proof}
$\Cov(f(Z),f(-Z))=\E[f(Z)f(-Z)]-\mu^2=-\mu^2$, so $\rho=-\mu^2/\sigma^2$; $\rho\ge-1$ gives
$\mu^2\le\sigma^2$, and the rest is \eqref{eq:cv-antithetic-var} and $1/(1+\rho)$. For the
call, $S_T(z)\,S_T(-z)=S^2e^{2\nu T}$. If $K\ge Se^{\nu T}$ and $S_T(z)>K$, then
$S_T(-z)<S^2e^{2\nu T}/K\le K$, so the two payoffs are never positive together. If
$K<Se^{\nu T}$, both are positive on $\{K<S_T<S^2e^{2\nu T}/K\}$, a nonempty interval of
values of $S_T$ and therefore an event of positive probability. The digital has the same
supports, and for puts the inequalities reverse.
\end{proof}

In the test case the median is $Se^{\nu T}=97.04<K$, so the at-the-money call is disjoint:
with $\mu=7.36429$ and $\sigma^2=158.006$, $\mu^2/\sigma^2=0.3432$, so $\rho=-0.3432$, the
standard-error ratio is $0.5730$ and the efficiency at $c=2$ is $1.523$, the course's
``about 57\%'' and ``about 1.5''. For strikes below the median the supports overlap,
$\E[f(Z)f(-Z)]>0$, and $\rho$ is larger than $-\mu^2/\sigma^2$ (\cref{fig:cv-antithetic}).

\begin{figure}[ht]
  \centering
  \includegraphics[width=\textwidth]{ch08-antithetic}
  \caption{(a) The discounted at-the-money call payoff $f(z)$, its antithetic image
    $f(-z)$ and the even part $f_{\mathrm e}$, the only part the antithetic estimator has to
    sample; the supports are disjoint (\cref{prop:cv-disjoint}). (b) The antithetic
    correlation $\rho$ against the strike, exact (lines, by quadrature) and from $20{,}000$
    pairs (dots); the dotted line is the median of $S_T$. Near the median the straddle
    $|S_T-K|$ is almost even in $z$: $\rho=0.85$ at the median and $0.90$ at its maximum,
    near $K=99$.}
  \label{fig:cv-antithetic}
\end{figure}

\begin{remark}[When pairing fails]\label{rem:cv-antithetic-fails}
Without monotonicity, \cref{prop:cv-antithetic}(iii) gives no guarantee. If $f$ is even,
then $f_{\mathrm o}=0$, $\rho=1$, and a pair is worth one draw at up to twice the cost. Payoffs
that are nearly even in $z$ (straddles, butterflies, digital ranges centred near the median of
$S_T$) have $\rho$ close to $1$ (\exref{ex:cv-straddle}). At the other extreme, a deep
out-of-the-money call is disjoint with $\mu^2\ll\sigma^2$, so $\rho\approx0$ and the
efficiency tends to $1$ at $c=2$: rare events need importance sampling
(\cref{ch:09-importance-sampling}). For paths, an arithmetic Asian call is nondecreasing in
every increment, so (iii) applies; a knock-out barrier is not monotone, and the
correlation must be measured.
\end{remark}

\demo{08-antithetic-control.ipynb}{the even and odd parts of any payoff, and the
antithetic correlation against the strike for calls, puts, digitals and straddles.}

\section{Moment matching}\label{sec:cv-mm}

Antithetic pairs make the sample symmetric. Moment matching goes further and transforms the
whole sample so that its first moments are exactly those of the model.

\begin{definition}[Moment matching]\label{def:cv-mm}
Let $Z_1,\dots,Z_n$ be i.i.d.\ $\mathcal N(0,1)$, $\bar Z_n=\frac1n\sum_iZ_i$ and
$\tilde s_n^2=\frac1n\sum_i(Z_i-\bar Z_n)^2$ (divisor $n$).
\begin{enumerate}[label=(\alph*)]
  \item \emph{Log matching}: $Z'_i=(Z_i-\bar Z_n)/\tilde s_n$. Then
    $\frac1n\sum_iZ'_i=0$ and $\frac1n\sum_iZ'^2_i=1$, so the simulated log-prices
    $\ln S+\nu T+sZ'_i$ have exactly the mean and the variance of $\ln S_T$.
  \item \emph{Exponential-mean matching}: $Z'_i=\alpha_n+Z_i/\tilde s_n$ with
    $\alpha_n=\frac s2-\frac1s\ln\big(\frac1n\sum_ie^{sZ_i/\tilde s_n}\big)$. The sample variance
    of the $Z'_i$ is still $1$, and $\frac1n\sum_ie^{sZ'_i}=e^{s^2/2}=\E e^{sZ}$: the sample
    mean of $S'_{T,i}$ equals the forward $Se^{(r-y)T}$.
\end{enumerate}
The moment-matched estimator is $\hat\mu^{\mathrm{mm}}_n=\frac1n\sum_if(Z'_i)$.
\end{definition}

In (b) the sample reprices the forward exactly, so the matched estimates of a call and a
put with the same strike satisfy put--call parity exactly. The course also matches the
``exponential variance'', which requires solving a nonlinear equation for the scale by
Newton's method. For paths, match each time step's normals across the sample.

\begin{proposition}[Moment matching: structure and bias]\label{prop:cv-moment-matching}
In the setting of \cref{def:cv-mm}(a), with $n\ge3$:
\begin{enumerate}[label=(\roman*)]
  \item $\hat\mu^{\mathrm{mm}}_n=\E f(Z)$ exactly for every quadratic $f(z)=a+bz+cz^2$.
  \item The vector $Z'=(Z'_1,\dots,Z'_n)$ is independent of $(\bar Z_n,\tilde s_n)$. Its
    components are identically distributed but not independent ($\sum_iZ'_i=0$).
  \item If $|f(z)|\le Ce^{c|z|}$ for all $z$, then, with $\mathrm{He}_4(z)=z^4-6z^2+3$,
    \[
      \E\hat\mu^{\mathrm{mm}}_n=\E f(Z)-\frac{\E[f(Z)\mathrm{He}_4(Z)]}{4n}+O(n^{-2}) .
    \]
  \item The sample variance of $f(Z'_1),\dots,f(Z'_n)$ does not estimate
    $n\Var\hat\mu^{\mathrm{mm}}_n$ in general. If the $n$ draws are split into $m$ batches of
    size $k=n/m$, each matched separately, the $m$ batch estimates are i.i.d.\ with mean
    $\E\hat\mu^{\mathrm{mm}}_k=\mu+O(1/k)$. For fixed $k$ and $m\to\infty$, the central limit
    theorem for these i.i.d.\ estimates gives an asymptotically valid confidence interval for
    $\E\hat\mu^{\mathrm{mm}}_k$; for fixed $m$ and $k\to\infty$, \cref{prop:mc-batch}(iii)
    gives one for $\mu$ provided $\sqrt k(\hat\mu^{\mathrm{mm}}_k-\mu)$ is asymptotically
    normal, which \cref{rem:cv-mm-as-cv} indicates but we do not prove.
\end{enumerate}
\end{proposition}

\begin{proof}
(i) $\frac1n\sum_i(a+bZ'_i+cZ'^2_i)=a+c=\E[a+bZ+cZ^2]$.

(ii) Complete $u_0=n^{-1/2}(1,\dots,1)$ to an orthonormal basis $u_0,\dots,u_{n-1}$ of
$\R^n$. The coordinates $G_j=\langle Z,u_j\rangle$ are i.i.d.\ $\mathcal N(0,1)$, being an
orthogonal transformation of a standard normal vector. Now $\bar Z_n=G_0/\sqrt n$,
$Z-\bar Z_n(1,\dots,1)=\sum_{j\ge1}G_ju_j$, so with $G=(G_1,\dots,G_{n-1})$ we get
$\tilde s_n=\|G\|/\sqrt n$ and $Z'=\sqrt n\sum_{j\ge1}(G_j/\|G\|)u_j$. In polar coordinates the
density of $G$ is a function of $\|G\|$ times the uniform measure on the unit sphere, so
the direction $G/\|G\|$ is independent of the radius $\|G\|$; both are independent of
$G_0$. Hence $Z'$, a function of the direction, is independent of $(\bar Z_n,\tilde s_n)$.
Permuting the $Z_i$ permutes the $Z'_i$, so they are exchangeable.

(iii) By exchangeability $\E\hat\mu^{\mathrm{mm}}_n=\E f(Z'_1)$. By (ii),
$Z'_1=\sqrt n\langle G/\|G\|,w\rangle$ with $w=(\langle e_1,u_j\rangle)_{j\ge1}$ and
$\|w\|^2=1-\frac1n$; by rotation invariance of $G$ this has the law of
$\sqrt{n-1}\,G_1/\|G\|$. As $G_1^2\sim\chi^2_1$ and $\|G\|^2-G_1^2\sim\chi^2_{n-2}$ are
independent, $G_1^2/\|G\|^2$ has the $\mathrm{Beta}(\frac12,\frac{n-2}2)$ law, and $Z'_1$ has the
density $q_n(x)=\kappa_n\big(1-\frac{x^2}{n-1}\big)^{(n-4)/2}$ on $|x|<\sqrt{n-1}$, and $0$
outside. Write $q_n=\kappa_n\sqrt{2\pi}\,\varphi\,g_n$, with $\varphi$ the standard normal
density and $g_n(x)=(1-u)^{(n-4)/2}e^{x^2/2}$, $u=x^2/(n-1)$. From
$\ln(1-u)=-u-\frac{u^2}2-\theta(u)$ with $0\le\theta(u)\le u^3/(1-u)$,
\[
  \ln g_n(x)=\frac{3x^2}{2(n-1)}-\frac{(n-4)x^4}{4(n-1)^2}-\frac{n-4}{2}\theta(u) .
\]
For $|x|\le n^{1/4}$ and $n\ge7$ we have $u\le\frac12$, so $\ln g_n(x)=h(x)/n+O((1+x^6)/n^2)$
with $h(x)=\frac32x^2-\frac14x^4$, and $|h|/n\le1$; hence
$g_n(x)=1+h(x)/n+O((1+x^{12})/n^2)$ uniformly there. For $|x|>n^{1/4}$, both $\varphi g_n\le
e^{-x^2(n-4)/(2(n-1))}\le e^{-x^2/4}$ and $\varphi(1+|h|/n)$ are dominated by a multiple of
$e^{-x^2/4}$, and against $|f|\le Ce^{c|x|}$ these regions contribute $O(e^{-\sqrt n/8})$ once
$n^{1/4}\ge8c$. Therefore $J_n(f):=\int f\varphi g_n\,\dd x=\E f(Z)+\E[f(Z)h(Z)]/n+O(n^{-2})$.
Since $\E f(Z'_1)=J_n(f)/J_n(1)$ and $J_n(1)=1+\E h(Z)/n+O(n^{-2})=1+\frac{3}{4n}+O(n^{-2})$,
\[
  \E f(Z'_1)=\E f(Z)+\frac{\E\big[f(Z)\big(h(Z)-\frac34\big)\big]}{n}+O(n^{-2}),
  \qquad h-\tfrac34=-\tfrac14\mathrm{He}_4 .
\]

(iv) For $f(z)=z$ the estimator is exactly $0=\E Z$, while the sample variance of the
$Z'_i$ is $1$. Batches drawn from disjoint sets of i.i.d.\ normals are independent and
identically distributed, with the bias $O(1/k)$ of (iii) and finite variance (the $Z'_i$ of a
batch are bounded by $\sqrt k$); the rest is the central limit theorem and
\cref{prop:mc-batch}(iii).
\end{proof}

\begin{remark}[Moment matching is a control variate in disguise]\label{rem:cv-mm-as-cv}
Expanding $f(Z'_i)\approx f(Z_i)+f'(Z_i)(Z'_i-Z_i)$ with
$Z'_i-Z_i\approx-\bar Z_n-Z_i(\tilde s_n-1)$ and $\tilde s_n-1\approx\frac12(\frac1n\sum_jZ_j^2-1)$,
and using Stein's identities $\E f'(Z)=\E[Zf(Z)]=:a_1$ and
$\E[Zf'(Z)]=\E[(Z^2-1)f(Z)]=:a_2$, one finds
\[
  \hat\mu^{\mathrm{mm}}_n\approx\frac1n\sum_{i=1}^n\Big(f(Z_i)-a_1Z_i-\frac{a_2}{2}(Z_i^2-1)\Big),
\]
so that, heuristically, $\sqrt n(\hat\mu^{\mathrm{mm}}_n-\mu)$ is asymptotically normal with variance
$\Var\big(f(Z)-a_1Z-\frac{a_2}2(Z^2-1)\big)$. Since $Z$ and $Z^2-1$ are uncorrelated with
variances $1$ and $2$, the coefficients $a_1$ and $a_2/2$ are exactly the optimal ones of
\cref{thm:cv-multiple}: \emph{asymptotically, log matching is the optimal control-variate
estimator with the controls $Z$ and $Z^2$}, obtained without estimating any coefficient
(\exref{ex:cv-skewness} proves this for $f(z)=z^3$). The bias, by contrast, is exact to first
order by \cref{prop:cv-moment-matching}(iii): its constant $-\frac14\E[f(Z)\mathrm{He}_4(Z)]$ is
$-6$ for $z^4$ and $1.75$ for the call (\cref{fig:cv-bias}). By this linearisation the test-case
standard-error ratios are $0.0162$ for the stock and $0.188$ for the call, the course's ``about 1.6\%'' and ``about 20\%'' (and its ``11\% and 33\% of the
antithetic errors'' are $0.0162/0.140$ and $0.188/0.573$). Exponential-mean matching
reprices the stock exactly and gives $0.182$ for the call.
\end{remark}

\begin{remark}[Batches: how many?]\label{rem:cv-mm-batches}
By \cref{prop:cv-moment-matching}(iv), with $m$ batches the bias is about $c_fm/n$ while the
standard error is of order $n^{-1/2}$, so their ratio grows like $m/\sqrt n$; with too few
batches the error estimate itself is noisy ($m-1$ degrees of freedom). Keep $m$ well below
$\sqrt n$: with $n=2^{20}$ and $m=256$ the bias of the call estimate is about
$1.75\cdot256/2^{20}\approx4\times10^{-4}$, a fifth of its standard error $2.2\times10^{-3}$.
How coverage collapses when $m$ is too large is the subject of \exref{ex:cv-mm-coverage}.
\end{remark}

\begin{remark}[Antithetic variates against moment matching]\label{rem:cv-anti-vs-mm}
Antithetic sampling acts \emph{locally}: each draw is paired with its mirror image,
independently of the others, so the pair averages are i.i.d.\ and the usual standard error
applies. The resulting sample is symmetric, so its mean and all its odd moments are exactly
$0$, but its variance is not matched. Moment matching acts \emph{globally}: one affine map,
computed from the whole sample, fixes the first two moments and introduces dependence and
an $O(1/n)$ bias. The two can be combined, and the order matters: rescaling a symmetric
sample keeps it symmetric, so antithetic sampling followed by scale matching fixes the mean,
the variance and every odd moment; but exponential-mean matching shifts the sample and
destroys its symmetry. With (randomised) quasi-Monte Carlo point sets
(\cref{ch:11-quasi-monte-carlo}) the low moments are already integrated very accurately, and
matching adds bias for little gain.
\end{remark}

\section{Common random numbers}\label{sec:cv-crn}

Often the quantity of interest is a difference of two prices: the same option under two
models or two volatilities, two strikes, or a bumped Greek (\cref{ch:13-mc-greeks}). Each
price may be poorly estimated while their difference is estimated well, provided both use
the same random numbers.

\begin{proposition}[Common random numbers]\label{prop:cv-crn}
Let $U$ be a random vector with independent components, $f_1(U),f_2(U)\in L^2$ with
variances $\sigma_1^2,\sigma_2^2$, and $\delta=\E f_1(U)-\E f_2(U)$.
\begin{enumerate}[label=(\roman*)]
  \item With independent draws for the two terms, the estimator of $\delta$ from $n$ draws
    of each has variance $(\sigma_1^2+\sigma_2^2)/n$; with common random numbers,
    $\frac1n\sum_i\big(f_1(U_i)-f_2(U_i)\big)$ has variance
    $\big(\sigma_1^2+\sigma_2^2-2\Cov(f_1(U),f_2(U))\big)/n$.
  \item If, for each $j$, $f_1$ and $f_2$ are monotone in $u_j$ in the same direction, then
    $\Cov(f_1(U),f_2(U))\ge0$ and common random numbers never increase the variance. If they
    are monotone in opposite directions in every coordinate, the covariance is $\le0$ and
    common random numbers never decrease it; if moreover $U\eqd-U$, pairing $f_1(U_i)$ with
    $f_2(-U_i)$ restores a nonnegative covariance without changing the mean.
  \item Let $f(\theta,u)$ depend on a parameter, with
    $|f(\theta+h,u)-f(\theta,u)|\le L(u)|h|$ and $\E L(U)^2<\infty$. Then the common-random-number
    difference quotient $(f(\theta+h,U)-f(\theta,U))/h$ has variance at most $\E L(U)^2$
    for every $h$, while with independent draws its variance
    $(\sigma^2(\theta+h)+\sigma^2(\theta))/h^2$ tends to infinity as $h\to0$ whenever
    $\sigma^2(\theta)=\Var f(\theta,U)>0$.
\end{enumerate}
\end{proposition}

\begin{proof}
(i) is the variance of a difference. (ii) As in the proof of
\cref{prop:cv-antithetic}(iii), let $V_j=\varepsilon_jU_j$ with signs chosen so that
$g_k(v)=f_k(\varepsilon_1v_1,\dots,\varepsilon_dv_d)$ is nondecreasing in every coordinate for
$k=1$; $V$ has independent components and $f_k(U)=g_k(V)$. In the first case $g_2$ is
nondecreasing as well, in the second $-g_2$ is, and \cref{lem:cv-association} gives
$\Cov(g_1(V),g_2(V))\ge0$, respectively $\le0$. For the pairing, $u\mapsto f_2(-u)$ moves in
the same direction as $f_1$ in every coordinate, so the first case applies to $f_1$ and
$f_2(-\,\cdot)$, and $f_2(-U)\eqd f_2(U)$. (iii) The variance is at most the second moment,
which is at most $\E L(U)^2$; the independent case is (i) divided by $h^2$.
\end{proof}

For the test case, the prices at $\sigma=20\%$ and $21\%$ differ by $0.39048$. Both
discounted payoffs are nondecreasing in $z$, and their correlation is $0.99999$: the
variance per draw of the difference is $0.535$ with common normals and $334.85$ with
independent ones, a standard-error ratio of $0.040$. Comparing a call with a put is the
opposite case (\exref{ex:cv-crn-parity}). Finite-difference Greeks are the main
application; \cref{prop:gr-finite-difference} works out the optimal bump with and without
common random numbers. \textcite{glasserman-yao-crn} give general conditions under which
common random numbers are guaranteed to help.

\section{Control variates}\label{sec:cv-control}

Suppose that alongside the payoff $Y=f(Z)$ we can compute, at little extra cost, a
\emph{control} $X=g(Z)$ whose mean $\mu_X$ is known exactly: the stock, a vanilla option,
a geometric average, a hedge. With i.i.d.\ copies $(X_i,Y_i)$ and a coefficient $\beta$,
\begin{equation}\label{eq:cv-estimator}
  \hat\mu_n(\beta)=\frac1n\sum_{i=1}^n\big(Y_i-\beta(X_i-\mu_X)\big).
\end{equation}
We estimate the mean of $Y$ by estimating the mean of the error $Y-\beta X$ and adding
back the known part. The course uses $\beta=1$.

\begin{theorem}[The optimal control coefficient]\label{thm:cv-optimal}
Let $X,Y\in L^2$ with $\sigma_X^2=\Var X>0$, $\sigma_Y^2=\Var Y$, $\sigma_{XY}=\Cov(X,Y)$ and
$\rho=\Corr(X,Y)$.
\begin{enumerate}[label=(\roman*)]
  \item For every $\beta$, $\hat\mu_n(\beta)$ is unbiased with
    $n\Var\hat\mu_n(\beta)=\sigma_Y^2-2\beta\sigma_{XY}+\beta^2\sigma_X^2$.
  \item This is minimised uniquely at $\beta^*=\sigma_{XY}/\sigma_X^2$, where
    \[
      n\Var\hat\mu_n(\beta^*)=\sigma_Y^2\,(1-\rho^2).
    \]
  \item If a sample with its control costs $c$ times a plain sample, the efficiency of
    $\hat\mu_n(\beta^*)$ relative to plain Monte Carlo is $1/\big(c(1-\rho^2)\big)$.
\end{enumerate}
\end{theorem}

\begin{proof}
(i) $\E[Y_i-\beta(X_i-\mu_X)]=\E Y$, and the variance of $Y-\beta X$ expands as stated.
(ii) Completing the square,
$\sigma_Y^2-2\beta\sigma_{XY}+\beta^2\sigma_X^2=\sigma_X^2(\beta-\beta^*)^2+\sigma_Y^2-\sigma_{XY}^2/\sigma_X^2$,
and $\sigma_{XY}^2/\sigma_X^2=\rho^2\sigma_Y^2$. (iii) is \eqref{eq:cv-efficiency}.
\end{proof}

Geometrically, $\beta^*(X-\mu_X)$ is the orthogonal projection of $Y-\E Y$ onto $X-\mu_X$ in
$L^2$, and the control removes all of $Y$'s variance except the part orthogonal to $X$. The
factor $1-\rho^2$ is unforgiving: $\rho=0.9$ gives a variance factor $0.19$, $\rho=0.99$ gives
$0.02$. With $\beta=1$ and $\sigma_X=\sigma_Y$ the variance factor is $2(1-\rho)$, the
course's formula; \exref{ex:cv-fixed-beta} asks when a fixed coefficient helps at all.

\begin{example}[The square-root call]\label{eg:cv-sqrt}
Take $Y=e^{-rT}(\sqrt{S_T}-\sqrt K)^+$. Since $\sqrt{S_T}$ is lognormal with volatility
$\sigma/2$, the ``pre-wash'' map of \cref{prop:sr-power} prices it as a Black--Scholes call
on $\sqrt S$ with strike $\sqrt K$, volatility $\sigma/2$ and yield
$\frac12(r+y+\sigma^2/4)$. The course chooses devious inputs, $S=K=10^4$, $T=1$, $r=1\%$,
$y=-1\%$, $\sigma=40\%$, which map exactly to the test case: the price is $7.36429$. The
control is the tangent of the payoff at the kink,
$X=e^{-rT}(S_T-K)^+/(2\sqrt K)$, that is $\frac1{200}$ of a vanilla call, with known mean
$\mu_X=8.43615$, about $15\%$ above the price (\cref{fig:cv-control}a). By quadrature,
$\rho=0.9960$, $\sigma_X^2=235.75$, $\sigma_Y^2=158.01$ and $\beta^*=0.8154$. With the
course's $\beta=1$ the standard-error ratio is $0.2424$, the course's ``errors reduced by
about 75\%''; with $\beta^*$ it is $\sqrt{1-\rho^2}=0.0889$, a further factor $7.4$ in
efficiency. The course inferred $\rho\approx0.97$ by reading the $\beta=1$ result through
$2(1-\rho)=0.25^2$, which assumes $\sigma_X=\sigma_Y$; here $\sigma_X/\sigma_Y=1.22$, and the
true correlation is much higher. (The course writes the resulting efficiency as $2/0.25^2=8$;
the formula is $1/(2\cdot0.25^2)=8$.)
\end{example}

\begin{figure}[ht]
  \centering
  \includegraphics[width=\textwidth]{ch08-control}
  \caption{The square-root call with the course's devious inputs. (a) The payoff
    $(\sqrt{S_T}-\sqrt K)^+$ and its control, the tangent call $(S_T-K)^+/(2\sqrt K)$.
    (b) Two thousand discounted samples $(X_i,Y_i)$ and the lines through the exact means
    $(\mu_X,\mu_Y)$ with slope $\beta=1$ and with the least-squares slope $\hat\beta$. Here $Y$ is a
    function of $X$, so the points lie on a curve; what the control leaves is the deviation of
    that curve from a straight line.}
  \label{fig:cv-control}
\end{figure}

\subsection{Estimating the coefficient}

$\beta^*$ involves $\Cov(X,Y)$, which is unknown (if it were known, we would hardly need
simulation). The natural estimator is the least-squares slope from the same sample,
\[
  \hat\beta_n=\frac{\sum_i(X_i-\bar X_n)(Y_i-\bar Y_n)}{\sum_i(X_i-\bar X_n)^2},\qquad
  \hat\mu^{\mathrm{cv}}_n=\bar Y_n-\hat\beta_n(\bar X_n-\mu_X),
\]
when the denominator is positive; when it vanishes (all $X_i$ equal) we set $\hat\beta_n=0$.
This is not an idle case: the tangent call of \cref{eg:cv-sqrt} vanishes with probability
$N(0.15)=0.56$, so all $n$ controls are $0$ with probability $0.56^n$.

\begin{proposition}[Estimated coefficient: bias and splitting]\label{prop:cv-estimated-beta}
Let $(X_i,Y_i)$ be i.i.d.\ copies of $(X,Y)$ as in \cref{thm:cv-optimal}.
\begin{enumerate}[label=(\roman*)]
  \item $\hat\beta_n\to\beta^*$ almost surely, and
    $\sqrt n\,(\hat\mu^{\mathrm{cv}}_n-\E Y)\Rightarrow\mathcal N\big(0,\sigma_Y^2(1-\rho^2)\big)$:
    estimating the coefficient costs nothing asymptotically.
  \item If $\E\hat\beta_n^2<\infty$, the bias is
    $\E\hat\mu^{\mathrm{cv}}_n-\E Y=-\Cov(\hat\beta_n,\bar X_n)$, and
    $|\E\hat\mu^{\mathrm{cv}}_n-\E Y|\le\sigma_X\,\mathrm{sd}(\hat\beta_n)/\sqrt n$; if moreover
    $\mathrm{sd}(\hat\beta_n)=O(n^{-1/2})$, the bias is $O(1/n)$.
  \item If $X$ has a continuous law and $\E[Y\mid X]=a+bX$ is affine, then
    $\E[\hat\mu^{\mathrm{cv}}_n\mid X_1,\dots,X_n]=\E Y$; so $\hat\mu^{\mathrm{cv}}_n$ is unbiased
    whenever it is integrable.
  \item (\emph{Splitting}.) Let $\hat\beta^{(1)}$ and $\hat\beta^{(2)}$ be the slopes computed
    on the two halves of the sample, and assume they are integrable. The estimator that
    applies each half's slope to the other half,
    \[
      \hat\mu^{\mathrm{split}}_n=\tfrac12\Big[\bar Y^{(1)}-\hat\beta^{(2)}(\bar X^{(1)}-\mu_X)\Big]
      +\tfrac12\Big[\bar Y^{(2)}-\hat\beta^{(1)}(\bar X^{(2)}-\mu_X)\Big],
    \]
    is unbiased and has the same limit law as in (i).
\end{enumerate}
\end{proposition}

\begin{proof}
(i) By the strong law, $\frac1n\sum(X_i-\bar X_n)(Y_i-\bar Y_n)\to\sigma_{XY}$ and
$\frac1n\sum(X_i-\bar X_n)^2\to\sigma_X^2$ almost surely, so $\hat\beta_n\to\beta^*$. Write
\[
  \sqrt n(\hat\mu^{\mathrm{cv}}_n-\E Y)=\sqrt n\big(\hat\mu_n(\beta^*)-\E Y\big)
  -(\hat\beta_n-\beta^*)\sqrt n(\bar X_n-\mu_X).
\]
The first term converges in law to $\mathcal N(0,\sigma_Y^2(1-\rho^2))$ by the central limit
theorem and \cref{thm:cv-optimal}; in the second, $\sqrt n(\bar X_n-\mu_X)$ converges in law
and $\hat\beta_n-\beta^*\to0$, so the product tends to $0$ in probability and Slutsky's lemma
\citep[Lemma~2.8]{vandervaart-asymptotic} concludes.

(ii) $\E\hat\mu^{\mathrm{cv}}_n=\E Y-\E[\hat\beta_n(\bar X_n-\mu_X)]$ and
$\E[\hat\beta_n(\bar X_n-\mu_X)]=\Cov(\hat\beta_n,\bar X_n)$ because $\E\bar X_n=\mu_X$.
Cauchy--Schwarz bounds the covariance by $\mathrm{sd}(\hat\beta_n)\,\sigma_X/\sqrt n$. Both
hypotheses can fail for small $n$: in the bivariate normal case
$\Var\hat\beta_n=\sigma_Y^2(1-\rho^2)/(\sigma_X^2(n-3))$ for $n\ge4$, but
$\E\hat\beta_3^2=\infty$ when $|\rho|<1$ (\exref{ex:cv-lavenberg-welch}).

(iii) Write $Y_i=a+bX_i+\varepsilon_i$ with $\E[\varepsilon_i\mid X_i]=0$; by independence
across $i$, $\E[\varepsilon_i\mid X_1,\dots,X_n]=0$. Almost surely the $X_i$ are distinct, and
$\hat\beta_n=b+\sum_i(X_i-\bar X_n)\varepsilon_i/\sum_i(X_i-\bar X_n)^2$ and
\[
  \hat\mu^{\mathrm{cv}}_n=a+b\mu_X+\bar\varepsilon_n-(\bar X_n-\mu_X)\frac{\sum_i(X_i-\bar X_n)\varepsilon_i}{\sum_i(X_i-\bar X_n)^2},
\]
which is $\E Y=a+b\mu_X$ plus a linear combination of the $\varepsilon_i$ with
coefficients that are functions of $X_1,\dots,X_n$; its conditional mean is $\E Y$.

(iv) Given the second half, $\hat\beta^{(2)}$ is a constant and the first half is still an
i.i.d.\ sample independent of it, so the first bracket has conditional mean $\E Y$ (for any
fixed $\beta$, $\hat\mu(\beta)$ is unbiased); it is integrable since
$\E|\hat\beta^{(2)}(\bar X^{(1)}-\mu_X)|=\E|\hat\beta^{(2)}|\,\E|\bar X^{(1)}-\mu_X|$. The same
holds for the second bracket. The limit law follows from (i) applied to each half, since
each half's slope converges to $\beta^*$.
\end{proof}

A first-order expansion of the ratio $\hat\beta_n$ suggests the size of the bias in (ii):
heuristically, $n(\E\hat\mu^{\mathrm{cv}}_n-\E Y)\to-\E\big[(X-\mu_X)^2\varepsilon\big]/\sigma_X^2$
with $\varepsilon=Y-\E Y-\beta^*(X-\mu_X)$, which vanishes in case (iii); \textcite{nelson-cv}
discusses this bias and its remedies. For the call with the
stock as control the constant is $-8.18$ (\cref{fig:cv-bias}). The bias is a fraction
$O(n^{-1/2})$ of the standard error and is negligible for large samples, but it matters
for small batches (for instance in the regressions inside American option pricing,
\cref{ch:14-american-mc}). The standard error of $\hat\mu^{\mathrm{cv}}_n$ is estimated by
the residual standard deviation with $n-2$ degrees of freedom. A pilot sample, used only to
estimate $\beta$ and then discarded, removes the bias like splitting does.
The exact small-sample variance in the normal case is computed in \exref{ex:cv-lavenberg-welch}.

\begin{figure}[ht]
  \centering
  \includegraphics[width=0.82\textwidth]{ch08-bias}
  \caption{$n$ times the bias of three estimators of the test-case call, from $6.4$ million
    normal draws at each sample size $n$ (bars: two standard errors). Moment matching
    (\cref{prop:cv-moment-matching}) and the control variate with in-sample coefficient
    (\cref{prop:cv-estimated-beta}) have bias of order $1/n$, with the predicted constants
    $1.75$ and $-8.18$ (dashed); the split-sample estimator is unbiased.}
  \label{fig:cv-bias}
\end{figure}

\subsection{Several controls}

With several controls $X=(X^1,\dots,X^q)$ of known mean $\mu_X\in\R^q$, use
$\hat\mu_n(\beta)=\frac1n\sum_i\big(Y_i-\beta^\top(X_i-\mu_X)\big)$ with $\beta\in\R^q$.

\begin{theorem}[Multiple controls as regression]\label{thm:cv-multiple}
Let $Y\in L^2$ with $\sigma_Y^2=\Var Y>0$, $X\in L^2(\R^q)$, with $\Sigma_{XX}=\Cov(X)$ positive definite and
$\Sigma_{XY}=\Cov(X,Y)\in\R^q$. Set
$\beta^*=\Sigma_{XX}^{-1}\Sigma_{XY}$ and $R^2=\Sigma_{XY}^\top\Sigma_{XX}^{-1}\Sigma_{XY}/\sigma_Y^2$.
\begin{enumerate}[label=(\roman*)]
  \item For every $\beta\in\R^q$,
    \[
      n\Var\hat\mu_n(\beta)=\sigma_Y^2(1-R^2)+(\beta-\beta^*)^\top\Sigma_{XX}(\beta-\beta^*),
    \]
    so $\beta^*$ is the unique optimum, with variance factor $1-R^2$, and $0\le R^2\le1$.
  \item $R^2=\max_{a\in\R^q\setminus\{0\}}\Corr(Y,a^\top X)^2$. If $\Sigma_{XY}\ne0$ the maximum
    is attained exactly at the nonzero multiples of $\beta^*$; if $\Sigma_{XY}=0$, every $a$ gives
    $0=R^2$. In particular $R^2\ge\Corr(Y,X^k)^2$ for each $k$, and adding a control never increases the
    optimal variance.
  \item If $S_{XX}=\sum_i(X_i-\bar X_n)(X_i-\bar X_n)^\top$ is invertible and
    $\hat\beta_n=S_{XX}^{-1}\sum_i(X_i-\bar X_n)(Y_i-\bar Y_n)$, then
    $\bar Y_n-\hat\beta_n^\top(\bar X_n-\mu_X)$ is the intercept $\hat a$ of the
    least-squares fit of $Y_i$ on $(1,X_i-\mu_X)$, that is, the fitted regression evaluated
    at $X=\mu_X$. As $n\to\infty$,
    $\sqrt n(\hat a-\E Y)\Rightarrow\mathcal N\big(0,\sigma_Y^2(1-R^2)\big)$.
\end{enumerate}
\end{theorem}

\begin{proof}
(i) $\Var(Y-\beta^\top X)=\sigma_Y^2-2\beta^\top\Sigma_{XY}+\beta^\top\Sigma_{XX}\beta$, and
completing the square with $\Sigma_{XY}=\Sigma_{XX}\beta^*$ gives the formula, since
$\beta^{*\top}\Sigma_{XX}\beta^*=\sigma_Y^2R^2$. The quadratic form is positive definite, so
the optimum is unique; at $\beta^*$ the variance $\sigma_Y^2(1-R^2)$ is nonnegative, and
$R^2\ge0$ since $\Sigma_{XX}^{-1}$ is positive definite.

(ii) In the inner product $\langle u,v\rangle_\Sigma=u^\top\Sigma_{XX}v$,
$a^\top\Sigma_{XY}=\langle a,\beta^*\rangle_\Sigma$, so by Cauchy--Schwarz
\[
  \Corr(Y,a^\top X)^2=\frac{(a^\top\Sigma_{XY})^2}{\sigma_Y^2\,a^\top\Sigma_{XX}a}
  \le\frac{\langle\beta^*,\beta^*\rangle_\Sigma}{\sigma_Y^2}=R^2,
\]
with equality if and only if $a$ and $\beta^*$ are linearly dependent, that is, when
$\beta^*\ne0$ (equivalently $\Sigma_{XY}\ne0$), if and only if $a$ is a multiple of $\beta^*$; if
$\Sigma_{XY}=0$ both sides vanish for every $a$. A larger set of controls maximises over a larger
set of combinations.

(iii) For fixed $b$, $\sum_i(Y_i-a-b^\top(X_i-\mu_X))^2$ is minimised by
$a=\bar Y_n-b^\top(\bar X_n-\mu_X)$; substituting leaves
$\sum_i\big((Y_i-\bar Y_n)-b^\top(X_i-\bar X_n)\big)^2$, minimised by $b=\hat\beta_n$ (normal
equations). The limit law follows as in \cref{prop:cv-estimated-beta}(i), with the strong law
for $S_{XX}/n$ and the multivariate central limit theorem.
\end{proof}

In terms of correlations, $1-R^2=1-\rho_{YX}^\top\rho_{XX}^{-1}\rho_{XY}$, the form used in the
course. Two cautions. First, each estimated coefficient costs a little: when $(X,Y)$ is
jointly normal, the variance of $\hat a$ is $\sigma_Y^2(1-R^2)/n$ times the loss factor
$(n-2)/(n-q-2)$ \citep{lavenberg-moeller-welch}, so $q$ must stay small compared with $n$.
Second, nearly collinear controls make $S_{XX}$ ill-conditioned and $\hat\beta_n$ erratic, even
though the fitted intercept may remain accurate; drop redundant controls. For the
square-root call, adding the stock to the tangent call raises $R^2$ from $0.9921$ to $0.9934$
and lowers the standard-error ratio from $0.0889$ to $0.0815$. The arithmetic Asian option with
geometric, vanilla and strip controls (\cref{prop:as-control}) is the course's showcase of
multiple controls.

\subsection{Nonlinear controls and hedges}

Two further ways to use known means lead back to linear controls.

\begin{proposition}[Nonlinear controls]\label{prop:cv-nonlinear}
\begin{enumerate}[label=(\roman*)]
  \item (\emph{Delta method}.) Let $h$ be defined and continuously differentiable on a
    neighbourhood $U$ of $(\mu_X,\E Y)$, with $h(\mu_X,y)=y$ whenever $(\mu_X,y)\in U$, and set
    $h(\bar X_n,\bar Y_n)=\bar Y_n$ when $(\bar X_n,\bar Y_n)\notin U$. Then $h(\bar X_n,\bar Y_n)\to\E Y$ almost
    surely and
    \[
      \sqrt n\big(h(\bar X_n,\bar Y_n)-\E Y\big)\Rightarrow
      \mathcal N\Big(0,\Var\big(Y-\beta_h(X-\mu_X)\big)\Big),\qquad
      \beta_h=-\partial_xh(\mu_X,\E Y):
    \]
    a nonlinear control is asymptotically the linear control with coefficient $\beta_h$, and
    no better than $\beta^*$.
  \item (\emph{Hedges as controls}.) Let $V(t,s)$ be the Black--Scholes price of a European
    claim with a Lipschitz payoff $F(S_T)$ (calls and puts qualify). For a grid
    $0=t_0<\dots<t_m=T$ let $\tilde S_t=e^{-(r-y)t}S_t$ and
    \[
      M_m=\sum_{k=0}^{m-1}e^{-yt_k}\,\partial_sV(t_k,S_{t_k})\,\big(\tilde S_{t_{k+1}}-\tilde S_{t_k}\big).
    \]
    Then $\E M_m=0$, and with $Y=e^{-rT}F(S_T)$,
    $\E\big[(Y-M_m-V(0,S))^2\big]\to0$ as $\max_k(t_{k+1}-t_k)\to0$.
\end{enumerate}
\end{proposition}

\begin{proof}
(i) By the strong law, $(\bar X_n,\bar Y_n)\in U$ for all large $n$, almost surely, so the
convention outside $U$ does not affect the limits; continuity of $h$ gives the almost sure
limit. By the multivariate
central limit theorem and the delta method \citep[Theorem~3.1]{vandervaart-asymptotic},
$\sqrt n(h(\bar X_n,\bar Y_n)-h(\mu_X,\E Y))$ converges in law to
$\nabla h^\top N$ with $N\sim\mathcal N(0,\Cov(X,Y))$, that is, to the law of
$\partial_xh\,X+\partial_yh\,Y$ centred. Differentiating $h(\mu_X,y)=y$ gives $\partial_yh=1$,
and $\partial_xh=-\beta_h$. The last claim is \cref{thm:cv-optimal}(ii).

(ii) Under $\Q$, $\tilde S_t=S\exp(\sigma W_t-\sigma^2t/2)$ is a square-integrable martingale and
the weights $e^{-yt_k}\partial_sV(t_k,S_{t_k})$ are bounded and $\F_{t_k}$-measurable, so every
term of $M_m$ has mean zero. Writing
$V(t,s)=e^{-r(T-t)}\E F\big(se^{\nu(T-t)+\sigma\sqrt{T-t}Z}\big)$ and differentiating under the
Gaussian integral, $V$ is $C^{1,2}$ on $[0,T)\times(0,\infty)$ and solves the Black--Scholes
equation; if $L$ is the Lipschitz constant of $F$, $|\partial_sV|\le Le^{-y(T-t)}$, the delta is
continuous, and $V(t,s)\to F(s)$ as $t\uparrow T$ by dominated convergence. It\^o's formula gives,
as in the delta-hedging argument of the IEOR 222 notes,
$\dd(e^{-rt}V(t,S_t))=e^{-rt}\partial_sV\,\sigma S_t\,\dd W_t=h_t\,\dd\tilde S_t$ with
$h_t=e^{-yt}\partial_sV(t,S_t)$. Letting $t\uparrow T$ (the integrand is bounded, and
$V(t,S_t)\to F(S_T)$ since $V$ is jointly continuous up to $T$ by the same argument), $Y=V(0,S)+\int_0^Th_t\,\dd\tilde S_t$. Hence
$Y-V(0,S)-M_m=\int_0^T(h_t-h_{\eta(t)})\,\dd\tilde S_t$ with $\eta(t)=t_k$ on $[t_k,t_{k+1})$,
and by the It\^o isometry
\[
  \E\big[(Y-M_m-V(0,S))^2\big]=\E\int_0^T(h_t-h_{\eta(t)})^2\sigma^2\tilde S_t^2\,\dd t\longrightarrow0
\]
by dominated convergence: the integrand is bounded by $4\sup|h|^2\sigma^2\tilde S_t^2$, which
is integrable, and tends to $0$ for every $t<T$ along every path, by continuity of $h$.
\end{proof}

Part (i) explains the classical ratio estimator $\bar Y_n\mu_X/\bar X_n$ ($\mu_X\ne0$, with
$U$ a neighbourhood not meeting $\{x=0\}$), for which
$\beta_h=\E Y/\mu_X$: for the call with the stock as control this is $0.075$, far from
$\beta^*=0.560$, so the ratio estimator barely helps. Part (ii) says that the ideal control is
the claim's own hedge: a control variate is a hedge whose cost is known. With the
Black--Scholes delta of the call itself as weights, the standard-error ratio for the test-case
call is about $0.48$, $0.15$ and $0.075$ for $m=1$, $12$ and $52$ rebalancing dates
($20{,}000$ paths), a variance falling asymptotically like $1/m$ as in \cref{prop:hd-discrete-error}. In
practice the delta of a nearby tractable claim is used: the static tangent call of
\cref{eg:cv-sqrt} is the one-date version, and the asymptotic square-root call of
\exref{ex:cv-asc} uses the vanilla call with the coefficient $F'(K+)=1$ read off the
replication.

\section{Combining the techniques: a laboratory}\label{sec:cv-lab}

\Cref{tab:cv-lab} runs every method of this chapter on the same $n=2^{20}$ normal draws
(seed $1$) and sets the measured standard-error ratios beside the exact asymptotic ones,
computed by quadrature from the formulas above. Moment matching uses $256$ separately matched
batches, and its measured ratios carry the noise of a batch standard error, about $4\%$.

\begin{table}[ht]
  \centering
  \small
  \begin{tabular}{@{}llccc@{}}
    \toprule
    target & method & exact & laboratory & efficiency \\
    \midrule
    stock & antithetic & $0.140$ & $0.140$ & $25.5$ \\
     & moment matching (log) & $0.016$ & $0.017$ & $1.9\times10^3$ \\
     & exponential matching; control $S_T$ & $0$ & $0$ & exact \\
    \midrule
    call & antithetic & $0.573$ & $0.574$ & $1.5$ \\
     & moment matching (log) & $0.188$ & $0.180$ & $14$ \\
     & moment matching (exponential) & $0.182$ & $0.174$ & $15$ \\
     & control $S_T$, $\hat\beta_n$ & $0.472$ & $0.472$ & $2.2$ \\
     & antithetic and control $S_T$ & $0.151$ & $0.152$ & $22$ \\
    \midrule
    square-root call & control, $\beta=1$ & $0.242$ & $0.242$ & $8.5$ \\
     & control, $\hat\beta_n$ & $0.089$ & $0.089$ & $63$ \\
     & tangent call and stock & $0.081$ & $0.082$ & $75$ \\
    \bottomrule
  \end{tabular}
  \caption{Standard-error ratios to plain Monte Carlo at equal $n$, exact (asymptotic, by
    quadrature) and measured in the laboratory ($n=2^{20}$, seed $1$), for the test case and
    the square-root call, and the efficiency $1/(2\varrho^2)$ of \eqref{eq:cv-efficiency} at
    the course's conservative cost $c=2$. Antithetic ratios compare $n$ pairs with $n$ draws.
    The exact entries for moment matching use the linearisation of \cref{rem:cv-mm-as-cv}.}
  \label{tab:cv-lab}
\end{table}

\begin{remark}[Why combinations are not multiplicative]\label{rem:cv-combine}
Antithetic sampling removes the odd part $f_{\mathrm o}$; a control applied afterwards
regresses the even part of $Y$ on the even part of $X$, and only the correlation of the even
parts matters. For the call with the stock as control, the even parts have correlation
$0.9645$, against $0.882$ for the raw payoffs, and the combination reaches $0.151$, better than
the product $0.573\times0.472=0.270$ of the separate ratios. The opposite also happens: the
control $Z$ itself, which helps alone, is useless after pairing, because its even part is
zero; the variance it would remove is odd, and pairing has already removed it. Likewise, log
matching already contains the optimal controls $Z$ and $Z^2$ (\cref{rem:cv-mm-as-cv}), so
little is left for a further control: exponential-mean matching, which also fixes the mean of
the stock, improves the ratio only from $0.188$ to $0.182$. The lesson of the table is the course's: almost any combination
helps, but the gains must be measured, not multiplied.
\end{remark}

\paragraph{Course examples reproduced.} The antithetic ratios $14\%$ (stock) and $57\%$ (call)
with efficiencies $25$ and $1.5$; the moment-matching ratios $1.6\%$ and $20\%$ (efficiencies
$1.9\times10^3$ and $14$ here, about $2000$ and $13$ in the course); the square-root call
$7.36429$, its control mean $8.43615$ and the $75\%$ error reduction at $\beta=1$; and the
call-payoff variance $161.2$.

\demo{08-antithetic-control.ipynb}{moment matching with batches, the bias of order $1/n$,
control variates with a slider for the coefficient, common random numbers for a volatility
bump, and the laboratory of \cref{tab:cv-lab} with measured execution times.}

\furtherreading
\textcite{glasserman-mc} is the standard reference: Sections~4.1 (control variates, with
multiple, small-sample and nonlinear controls in~4.1.2--4.1.4), 4.2 (antithetic variates),
4.5.1 (moment matching) and~7.1 (common random numbers for Greeks). Antithetic variates go back to
\textcite{hammersley-morton}; the association inequality behind
\cref{lem:cv-association} is developed in \textcite{esary-proschan-walkup}.
\textcite{lavenberg-welch} survey control variates with estimated coefficients; the loss
factor is due to \textcite{lavenberg-moeller-welch}, and \textcite{nelson-cv} treats the bias
and its remedies. \textcite{glasserman-yao-crn} study common random numbers.
\textcite{boyle-broadie-glasserman} review variance reduction for security pricing.
